\documentclass[11pt]{article}
\usepackage[margin=1in]{geometry}
\usepackage{amsmath,amssymb,amsthm}
\usepackage{hyperref}

\newtheorem{observation}{Observation}
\theoremstyle{definition}
\newtheorem{deadend}{Dead end}

\newcommand{\Z}{\mathbb{Z}}

\title{What fails, and why: negative results on the Umans--Wang \\
$(\alpha,\beta)$-divisor conjecture}
\author{Eric Schultz\\ \small independent researcher \\ \small ORCID 0009-0006-6283-1696}
\date{July 2026}

\begin{document}
\maketitle

\begin{abstract}
This companion note documents the approaches we tried against the Umans--Wang
$(\alpha,\beta)$-divisor conjecture that \emph{failed}, and---more usefully---why
each failed. The central empirical lesson is that every elementary ``barrier''
argument we produced was, on exact re-examination, an artifact: a
generic/average magnitude estimate mistaken for a lower bound, a degenerate
solution, or a solver timeout misread as infeasibility. We catalogue these so that
future work does not repeat them, and we map the construction families that do and
do not help, localizing the one mechanism (algebraic factoring of differences) that
gives genuine but only constant-factor benefit in $\Z$.
\end{abstract}

\section{Scope}
The main note \cite{main} proves $2\beta+\alpha\ge 1$ and reduces the conjecture at
fixed $\alpha$ to a lattice-clustering statement equivalent to the conjecture. Here
we record the negative space around that result: the arguments that looked like
progress and were not. All claims below were checked by exact computation; the
recurring theme is that exact checks repeatedly overturned plausible heuristics.

\section{The recurring artifact: generic bounds are not lower bounds}

Every ``barrier'' (i.e.\ would-be refutation of the conjecture) we found reduced to
the same error, in five guises.

\begin{deadend}[Vertex-log / incident-log forcing]
One argues: a coordinate $s_i$ constrained modulo a product of primes of total log
$L$ ``must'' have magnitude $\approx e^{L}$, and since some vertex accumulates
$L\gg n^{\alpha}$, the budget is exceeded. \textbf{Why it fails:} this is the
Minkowski \emph{average}, not a lower bound. Residue sharing---choosing neighbours'
values so the forced residue is small---defeats it. Exact search finds solutions
below the average in every tested case.
\end{deadend}

\begin{deadend}[Row-product forcing in the asymmetric ($S-T$) model]
The two-sided analogue: $s_i$ is ``forced'' modulo the product of its row's primes
$\approx e^{n/2K}$. \textbf{Why it fails:} identical to the above. With two free
sets, choosing the $t_j$ pulls $s_i$ strictly below the row product; exact search
confirms the gap grows with $K$.
\end{deadend}

\begin{deadend}[Degenerate covers]
Small ``solutions'' returned by unconstrained search often place $s_i=t_j$, so the
difference is $0$ and ``covers'' a prime trivially and invalidly. \textbf{Why it
fails:} a zero difference is not a cover; forbidding $s_i=t_j$ (and requiring
distinctness \emph{within} each set only) removes these. Some genuine solutions
survive, but the degenerate ones must be filtered first.
\end{deadend}

\begin{deadend}[Timeout as infeasibility]
An integer program for the minimum $K$ or minimum magnitude returns
``not solved'' (a solver time-out); this is read as ``infeasible,'' inflating the
apparent minimum and manufacturing a barrier. \textbf{Why it fails:} time-out is
not infeasibility. Distinguishing \texttt{Optimal}/\texttt{Infeasible}/\texttt{Timeout}
explicitly, and cross-checking small cases against brute force, dissolves the effect;
at the only cleanly solved sizes the counting floor is in fact reachable.
\end{deadend}

\begin{deadend}[Wrong magnitude regime]
Small-$n$ experiments (magnitude $\le 2n$, or $n^{O(1)}$) are used to estimate the
scaling exponent. \textbf{Why it fails:} the conjecture's difficulty lives in the
\emph{bundled} regime, where each difference is a product of $\Theta(n^{\alpha}/\log
n)$ primes $>n/2$ and magnitude is $e^{n^{\alpha}}$. A difference $\le 2n$ carries
at most one prime $>n/2$, so no polynomial-magnitude experiment ever enters the
regime that matters; observed exponents there are about the wrong quantity.
\end{deadend}

\begin{observation}
The uniform moral: in this problem, no counting or forcing argument yields a valid
magnitude lower bound, because the covering freedom (residue sharing across a
coprime system) always beats the generic estimate. The only genuine obstruction is
the lattice-clustering statement of \cite{main}, which is equivalent to the
conjecture itself. \emph{Any future ``barrier'' claim derived from a counting or
forcing bound should be presumed an artifact until checked exactly.}
\end{observation}

\section{Construction families: what helps and what does not}

On the constructive (pro-conjecture) side we tested many explicit families,
measuring how many primes of $(n/2,n]$ a typical difference bundles relative to a
random set.

\begin{itemize}
\item \textbf{Perfect squares: the one $\Z$ win.} With $S=\{a_i^2\}$, differences
factor as $a_i^2-a_j^2=(a_i-a_j)(a_i+a_j)$, giving two independent large factors and
$\approx 1.4$--$2\times$ the bundling of a random set. This is a genuine mechanism
but only a \emph{constant} factor; it does not change the exponent.
\item \textbf{Higher powers: no.} For $a^k-b^k$ with $k\ge 3$, staying within the
magnitude budget forces the roots below $B^{1/k}$, collapsing the extra factors;
bundling drops back to generic. The benefit peaks sharply at $k=2$.
\item \textbf{Factorials: structurally useless.} A prime $p>n/2$ divides $m!$ only
if $m\ge p$, but $m!\le B$ forces $m$ small; factorials are smooth in \emph{small}
primes only, the wrong range.
\item \textbf{Near-squares ($a\cdot b$, $a\approx b$): no.} Bundling is generic.
This sharpens the squares result: the benefit comes from the \emph{exact algebraic
identity} $a^2-b^2=(a-b)(a+b)$, not merely from two roughly equal factors.
\item \textbf{Fractal / digit-restricted / Beatty / Fibonacci: no.} These impose
\emph{additive} structure (small difference sets), which is the sum--product
antagonist of the multiplicative richness covering needs. The golden ratio, despite
appearing in the unit-separation lemma of \cite{BSSZ}, does not transfer: its role
there is archimedean separation, unrelated to integer divisibility.
\item \textbf{Naive $G\cdot P$ in $\Z$: no better than random.} A sparse geometric
$G$ is the known-inadequate Balog--Wooley example; a dense multiplicative $G$
(smooth numbers) matches random, because in the accessible regime random already
covers. The genuine $G\cdot P$ power requires the \emph{unit lattice} of a number
field, which has no faithful $\Z$ analogue.
\end{itemize}

\begin{observation}
Every family that helps does so through one mechanism: making differences factor
into several large independent pieces. In $\Z$ this tops out at two factors
(squares); reaching $\Theta(n^{\alpha}/\log n)$ factors per difference is inherently
multi-dimensional and lands in the number-field construction of \cite{BSSZ}. The
construction search is, in this sense, saturated: the ceiling in $\Z$ is understood,
and the only route past it is the one identified in the main note.
\end{observation}

\section{Computational limits}

Exact methods (bitmask minimum-$K$ search; brute-force covering; small integer
programs) reach roughly $n\le 34$ for minimum-$K$ and $K\le 3$--$8$ for the
congruence systems, before the exponential wall. The exact minimum-$K$ sequence
does show the gap to the counting floor staying in $\{0,1\}$ up to $n=34$---mildly
suggestive that the floor is reachable---but this is in the no-bundling regime and
therefore not evidence about the conjecture's true (bundled) regime. No computation
available to us reaches that regime, which requires super-polynomial magnitudes;
this is not a limitation of implementation but of the problem's location.

\section{Summary}
The negative results converge on a single picture consistent with the main note: the
elementary obstructions are mirages (five documented forms of one error), the
$\Z$-native constructions saturate at a constant-factor gain from the difference-of-
squares identity, and the true difficulty is the number-field clustering question.
Recording these dead ends is the practical contribution here: they are the specific
paths a future attack can skip.

\begin{thebibliography}{9}
\bibitem{main} E.~Schultz. \emph{A magnitude lower bound and a lattice-clustering
reformulation of the Umans--Wang $(\alpha,\beta)$-divisor conjecture.} Companion
paper, 2026.
\bibitem{UW} C.~Umans and S.~Wang. \emph{A number-theoretic conjecture implying
faster algorithms for polynomial factorization and integer factorization.}
arXiv:2511.10851, 2025.
\bibitem{BSSZ} T.~F.~Bloom, W.~Sawin, C.~Schildkraut, and D.~Zhelezov.
\emph{The sum-product conjecture is false for real numbers.} arXiv:2605.28781, 2026.
\end{thebibliography}

\end{document}
